\documentclass[aspectratio=169]{beamer}
% ===== shared Beamer preamble: light Metropolis, Libertinus serif =====
% Projection-first default: WHITE background, DARK (near-black) text. Dark-on-light
% is what survives a real projector; light-on-dark washes out in a lit hall.
% Font: Libertinus, an elegant old-style book serif whose matching Libertinus
% Math puts text and formulas in one voice. (A sans course may swap in lmodern
% with \usefonttheme[onlymath]{serif} instead; do NOT use a designer webfont.)
%
% Engine: pdflatex (no fontspec needed). Compile twice for the nav bar / toc.
%
% Each deck does:
%   \documentclass[aspectratio=169]{beamer}
%   % ===== shared Beamer preamble: light Metropolis, Libertinus serif =====
% Projection-first default: WHITE background, DARK (near-black) text. Dark-on-light
% is what survives a real projector; light-on-dark washes out in a lit hall.
% Font: Libertinus, an elegant old-style book serif whose matching Libertinus
% Math puts text and formulas in one voice. (A sans course may swap in lmodern
% with \usefonttheme[onlymath]{serif} instead; do NOT use a designer webfont.)
%
% Engine: pdflatex (no fontspec needed). Compile twice for the nav bar / toc.
%
% Each deck does:
%   \documentclass[aspectratio=169]{beamer}
%   % ===== shared Beamer preamble: light Metropolis, Libertinus serif =====
% Projection-first default: WHITE background, DARK (near-black) text. Dark-on-light
% is what survives a real projector; light-on-dark washes out in a lit hall.
% Font: Libertinus, an elegant old-style book serif whose matching Libertinus
% Math puts text and formulas in one voice. (A sans course may swap in lmodern
% with \usefonttheme[onlymath]{serif} instead; do NOT use a designer webfont.)
%
% Engine: pdflatex (no fontspec needed). Compile twice for the nav bar / toc.
%
% Each deck does:
%   \documentclass[aspectratio=169]{beamer}
%   \input{beamer-preamble}
% and is compiled FROM the directory that holds this file so \input resolves.
%
% House rule: ASCII punctuation only. No prose "---" (renders an em dash) and
% no prose "--" (renders an en dash). The ONLY allowed "--" is a TikZ line
% segment, e.g. \draw (a) -- (b). Use a comma, colon, parentheses, or "to".

\usepackage[T1]{fontenc}
\usepackage{libertinus}   % Libertinus Serif + Biolinum + matching Libertinus Math
\usefonttheme{serif}      % body, headings, and math all in the serif

\usetheme{metropolis}
\metroset{
  sectionpage=progressbar,
  numbering=fraction,
  progressbar=foot,
  block=fill
}

% --- accents for a LIGHT canvas (same names; retune hues per course) ---
\definecolor{accent}{HTML}{0369A1}    % deep sky blue (primary accent)
\definecolor{accent2}{HTML}{B45309}   % dark amber (secondary highlight)
\definecolor{good}{HTML}{047857}      % dark green (success / positive state)
\definecolor{bad}{HTML}{BE123C}       % dark rose (failure / negative state)
\colorlet{panel}{black!8}             % a panel slightly darker than the bg

% MANDATORY: white background, near-black text (the projection-first default).
\setbeamercolor{normal text}{fg=black!87, bg=white}
\setbeamercolor{alerted text}{fg=accent}
\setbeamercolor{progress bar in head/foot}{fg=accent}
\setbeamercolor{progress bar in section page}{fg=accent}
\setbeamercolor{title separator}{fg=accent}
% MANDATORY: frametitle as a pale panel strip with near-black text.
\setbeamercolor{frametitle}{fg=black!88, bg=panel}
\setbeamercolor{frametitle right}{fg=black!88, bg=panel}

% --- math: Libertinus already supplies it; amssymb/dsfont for extras ---
\usepackage{amsmath, amssymb}
\usepackage{dsfont}   % \mathds{1} indicator (amssymb has no blackboard-bold digit 1)

% --- TikZ, styled dark-on-light so every diagram reads on a projector ---
\usepackage{tikz}
\usetikzlibrary{arrows.meta, positioning, calc, backgrounds, fit, shapes.geometric}
\tikzset{>={Stealth[round]}}
\colorlet{fg}{black!85}
% MANDATORY: keep muted DARK enough on the light background. black!55 is the
% legible floor; anything lighter washes out under a projector. Do not raise it.
\colorlet{muted}{black!55}
\tikzset{
  every node/.style={text=fg, font=\footnotesize},
  edge/.style={draw=fg, thick, ->},
  dimedge/.style={draw=muted, thick, ->},
  nodebox/.style={draw=fg, thick, rounded corners=2pt, inner sep=5pt, text=fg},
  accentbox/.style={draw=accent, line width=1pt, rounded corners=2pt, inner sep=5pt, text=fg},
  gridcell/.style={draw=muted, minimum size=1.05cm, anchor=center, text=fg},
  treenode/.style={circle, draw=fg, thick, minimum size=0.9cm, inner sep=1pt, text=fg},
  good fill/.style={fill=good!14!white, draw=good, text=fg},
  bad fill/.style={fill=bad!12!white, draw=bad, text=fg},
  acc fill/.style={fill=accent!13!white, draw=accent, text=fg},
}

% --- minimal-text helper: one big centered takeaway line ---
\newcommand{\takeaway}[1]{\vfill\begin{center}\Large\alert{#1}\end{center}\vfill}

% and is compiled FROM the directory that holds this file so \input resolves.
%
% House rule: ASCII punctuation only. No prose "---" (renders an em dash) and
% no prose "--" (renders an en dash). The ONLY allowed "--" is a TikZ line
% segment, e.g. \draw (a) -- (b). Use a comma, colon, parentheses, or "to".

\usepackage[T1]{fontenc}
\usepackage{libertinus}   % Libertinus Serif + Biolinum + matching Libertinus Math
\usefonttheme{serif}      % body, headings, and math all in the serif

\usetheme{metropolis}
\metroset{
  sectionpage=progressbar,
  numbering=fraction,
  progressbar=foot,
  block=fill
}

% --- accents for a LIGHT canvas (same names; retune hues per course) ---
\definecolor{accent}{HTML}{0369A1}    % deep sky blue (primary accent)
\definecolor{accent2}{HTML}{B45309}   % dark amber (secondary highlight)
\definecolor{good}{HTML}{047857}      % dark green (success / positive state)
\definecolor{bad}{HTML}{BE123C}       % dark rose (failure / negative state)
\colorlet{panel}{black!8}             % a panel slightly darker than the bg

% MANDATORY: white background, near-black text (the projection-first default).
\setbeamercolor{normal text}{fg=black!87, bg=white}
\setbeamercolor{alerted text}{fg=accent}
\setbeamercolor{progress bar in head/foot}{fg=accent}
\setbeamercolor{progress bar in section page}{fg=accent}
\setbeamercolor{title separator}{fg=accent}
% MANDATORY: frametitle as a pale panel strip with near-black text.
\setbeamercolor{frametitle}{fg=black!88, bg=panel}
\setbeamercolor{frametitle right}{fg=black!88, bg=panel}

% --- math: Libertinus already supplies it; amssymb/dsfont for extras ---
\usepackage{amsmath, amssymb}
\usepackage{dsfont}   % \mathds{1} indicator (amssymb has no blackboard-bold digit 1)

% --- TikZ, styled dark-on-light so every diagram reads on a projector ---
\usepackage{tikz}
\usetikzlibrary{arrows.meta, positioning, calc, backgrounds, fit, shapes.geometric}
\tikzset{>={Stealth[round]}}
\colorlet{fg}{black!85}
% MANDATORY: keep muted DARK enough on the light background. black!55 is the
% legible floor; anything lighter washes out under a projector. Do not raise it.
\colorlet{muted}{black!55}
\tikzset{
  every node/.style={text=fg, font=\footnotesize},
  edge/.style={draw=fg, thick, ->},
  dimedge/.style={draw=muted, thick, ->},
  nodebox/.style={draw=fg, thick, rounded corners=2pt, inner sep=5pt, text=fg},
  accentbox/.style={draw=accent, line width=1pt, rounded corners=2pt, inner sep=5pt, text=fg},
  gridcell/.style={draw=muted, minimum size=1.05cm, anchor=center, text=fg},
  treenode/.style={circle, draw=fg, thick, minimum size=0.9cm, inner sep=1pt, text=fg},
  good fill/.style={fill=good!14!white, draw=good, text=fg},
  bad fill/.style={fill=bad!12!white, draw=bad, text=fg},
  acc fill/.style={fill=accent!13!white, draw=accent, text=fg},
}

% --- minimal-text helper: one big centered takeaway line ---
\newcommand{\takeaway}[1]{\vfill\begin{center}\Large\alert{#1}\end{center}\vfill}

% and is compiled FROM the directory that holds this file so \input resolves.
%
% House rule: ASCII punctuation only. No prose "---" (renders an em dash) and
% no prose "--" (renders an en dash). The ONLY allowed "--" is a TikZ line
% segment, e.g. \draw (a) -- (b). Use a comma, colon, parentheses, or "to".

\usepackage[T1]{fontenc}
\usepackage{libertinus}   % Libertinus Serif + Biolinum + matching Libertinus Math
\usefonttheme{serif}      % body, headings, and math all in the serif

\usetheme{metropolis}
\metroset{
  sectionpage=progressbar,
  numbering=fraction,
  progressbar=foot,
  block=fill
}

% --- accents for a LIGHT canvas (same names; retune hues per course) ---
\definecolor{accent}{HTML}{0369A1}    % deep sky blue (primary accent)
\definecolor{accent2}{HTML}{B45309}   % dark amber (secondary highlight)
\definecolor{good}{HTML}{047857}      % dark green (success / positive state)
\definecolor{bad}{HTML}{BE123C}       % dark rose (failure / negative state)
\colorlet{panel}{black!8}             % a panel slightly darker than the bg

% MANDATORY: white background, near-black text (the projection-first default).
\setbeamercolor{normal text}{fg=black!87, bg=white}
\setbeamercolor{alerted text}{fg=accent}
\setbeamercolor{progress bar in head/foot}{fg=accent}
\setbeamercolor{progress bar in section page}{fg=accent}
\setbeamercolor{title separator}{fg=accent}
% MANDATORY: frametitle as a pale panel strip with near-black text.
\setbeamercolor{frametitle}{fg=black!88, bg=panel}
\setbeamercolor{frametitle right}{fg=black!88, bg=panel}

% --- math: Libertinus already supplies it; amssymb/dsfont for extras ---
\usepackage{amsmath, amssymb}
\usepackage{dsfont}   % \mathds{1} indicator (amssymb has no blackboard-bold digit 1)

% --- TikZ, styled dark-on-light so every diagram reads on a projector ---
\usepackage{tikz}
\usetikzlibrary{arrows.meta, positioning, calc, backgrounds, fit, shapes.geometric}
\tikzset{>={Stealth[round]}}
\colorlet{fg}{black!85}
% MANDATORY: keep muted DARK enough on the light background. black!55 is the
% legible floor; anything lighter washes out under a projector. Do not raise it.
\colorlet{muted}{black!55}
\tikzset{
  every node/.style={text=fg, font=\footnotesize},
  edge/.style={draw=fg, thick, ->},
  dimedge/.style={draw=muted, thick, ->},
  nodebox/.style={draw=fg, thick, rounded corners=2pt, inner sep=5pt, text=fg},
  accentbox/.style={draw=accent, line width=1pt, rounded corners=2pt, inner sep=5pt, text=fg},
  gridcell/.style={draw=muted, minimum size=1.05cm, anchor=center, text=fg},
  treenode/.style={circle, draw=fg, thick, minimum size=0.9cm, inner sep=1pt, text=fg},
  good fill/.style={fill=good!14!white, draw=good, text=fg},
  bad fill/.style={fill=bad!12!white, draw=bad, text=fg},
  acc fill/.style={fill=accent!13!white, draw=accent, text=fg},
}

% --- minimal-text helper: one big centered takeaway line ---
\newcommand{\takeaway}[1]{\vfill\begin{center}\Large\alert{#1}\end{center}\vfill}


% ===== <DECK TITLE> animated slide deck (skeleton) =====
% Engine: LuaLaTeX. Compile FROM the slides/ directory so \input resolves:
%   lualatex <thisfile>.tex   (run twice for the nav/toc)
%
% House rules baked into this skeleton:
%   * MINIMAL TEXT: a few words per slide, no paragraphs, no bullet walls.
%   * A DIAGRAM on essentially every content slide (TikZ, styled light-on-dark).
%   * EVERY piece appears on its own click: [<+->] on lists, \pause between
%     blocks, \only<n->/\onslide<n-> on TikZ elements.
%   * AT MOST ONE big formula per slide, built up term by term with \only.
%   * Follow the section structure of the source lecture note.
%   * NO prose "---" or "--". The only allowed "--" is a TikZ "(a) -- (b)" segment.
%   * Every number is grounded in the source notes / spec; do not invent values.

\title{<DECK TITLE>}
\subtitle{<one-line subtitle, lower case is fine>}
\author{<Course / Weekend>}
\date{}

\begin{document}

% ------------------------------------------------------------------ 1. title
\maketitle

% ------------------------------------------------------------------ 2. a section per source-note section
\section{<First section of the source note>}

% PATTERN A: diagram first, one takeaway revealed on the second click.
\begin{frame}{<short frame title, sentence case>}
  \centering
  \begin{tikzpicture}[node distance=8mm and 16mm]
    \node[acc fill, nodebox] (a) {<thing>};
    \node[nodebox, right=of a] (b) {<thing>};
    % the edge is hidden until click 2 so the relation appears piece by piece
    \only<2->{\draw[edge] (a) -- (b);}
  \end{tikzpicture}

  \only<2->{\takeaway{<the one idea of this slide, a few words>}}
\end{frame}

% PATTERN B: [<+->] reveals each bullet on its own click; diagram alongside.
\begin{frame}{<short frame title>}
  \centering
  \begin{tikzpicture}[x=20mm]
    \foreach \i/\lab in {0/A,1/B,2/C} {
      \node[nodebox, minimum width=15mm, minimum height=18mm] (n\i) at (\i,0) {};
      \node at (\i,0.5) {\small \lab};
    }
    % a labelled detail that only appears on click 2, styled light-on-dark
    \only<2->{\node[text=muted] at (1,-0.05) {\small <detail>};}
  \end{tikzpicture}

  \begin{itemize}[<+->]
    \item <first point, a few words>
    \item \textbf{<contrast>}: <a few words>
    \item \textbf{<contrast>}: <a few words>
  \end{itemize}
\end{frame}

% ------------------------------------------------------------------ 3. one formula, built term by term
\section{<The formula section>}

% PATTERN C: at most ONE big formula, assembled one \only-term at a time.
% Each \only<n> REPLACES the previous so only one version shows per click.
% \underbrace names each term as it lands. The "%" after each \only suppresses
% a spurious space so the formula does not jump horizontally between clicks.
\begin{frame}{<the rule, in words>}
  \centering
  \only<1>{\Large $\mathrm{SCORE}(x) \;=\; \bar V(x)$}%
  \only<2>{\Large $\mathrm{SCORE}(x) \;=\; \underbrace{\bar V(x)}_{\text{exploit}}$}%
  \only<3->{\Large $\mathrm{SCORE}(x) \;=\; \underbrace{\bar V(x)}_{\text{exploit}}
      \;+\; \underbrace{w\sqrt{\dfrac{\ln N}{n(x)}}}_{\text{explore}}$}

  \vspace{4mm}
  \begin{itemize}[<+->]
    \item <what the first term does, a few words>
    \item <what the second term does, a few words>
    \item <the knob / the consequence>
  \end{itemize}

  \only<5->{\takeaway{<the punchline, a few words>}}
\end{frame}

% ------------------------------------------------------------------ 4. a multi-step diagram (phases)
\section{<The process section>}

% PATTERN D: ONE fixed diagram, each phase layered with \only on TikZ + a
% synchronized description list. Build each \item<n-> to match the \only<n>.
\begin{frame}{<process>, step by step}
  \centering
  \begin{tikzpicture}[node distance=9mm and 13mm]
    \node[treenode, acc fill] (r) {};
    \node[treenode, below left=of r] (l) {};
    \node[treenode, below right=of r] (m) {};
    \draw[edge] (r) -- (l);
    \draw[edge] (r) -- (m);

    % step 1: highlight a path
    \only<1>{\draw[edge, accent, line width=1.4pt] (r) -- (m);}
    \only<2->{\draw[edge] (r) -- (m);}
    % step 2: add a new node
    \only<2->{\node[treenode, good fill, below=9mm of m] (leaf) {};
              \draw[edge] (m) -- (leaf);}
    % step 3: a dashed action arrow (positioned so its label clears other nodes)
    \only<3>{\draw[edge, accent2, densely dashed] (leaf) -- ++(0.9,-1.1)
               node[right, text=accent2]{\scriptsize <action>};}
    % step 4: arrows climbing back
    \only<4>{\draw[edge, good, line width=1.2pt] (leaf) to[bend left=12] (m);
             \draw[edge, good, line width=1.2pt] (m) to[bend left=12] (r);}
  \end{tikzpicture}

  \medskip
  \begin{description}
    \item<1->[\alert{1 <step>}] <a few words>
    \item<2->[\alert{2 <step>}] <a few words>
    \item<3->[\alert{3 <step>}] <a few words>
    \item<4->[\alert{4 <step>}] <a few words>
  \end{description}
\end{frame}

% ------------------------------------------------------------------ 5. a worked example with REAL numbers
\section{<The worked example>}

% PATTERN E: two columns, the running example. Build the figure click by click
% with \only; keep every number identical to the source spec. Place node labels
% with care (see the dark-background collision rule in the skill) so a caption
% never lands on top of another label or runs through a node.
\begin{frame}{<the running example>}
  \begin{columns}[c]
    \column{0.46\textwidth}
      \centering
      \begin{tikzpicture}[x=11mm, y=11mm]
        \foreach \r in {0,1,2,3}{ \foreach \c in {0,1}{
          \node[gridcell, minimum size=0.95cm] at (\c, -\r) {};
        }}
        \only<2->{\node[text=accent, font=\bfseries] at (0,0) {<S>};}
        \only<3->{\node[good fill, gridcell, minimum size=0.95cm] at (0,-3) {};
                  \node[text=good, font=\bfseries] at (0,-3) {<T>};}
      \end{tikzpicture}
    \column{0.50\textwidth}
      \begin{itemize}[<+->]
        \item <label> $=(0,0)$
        \item <label> $=(3,0)$, reward $1.0$
      \end{itemize}
  \end{columns}
\end{frame}

% ------------------------------------------------------------------ 6. the bridge / takeaway
\section{<Bridge or punchline>}

\begin{frame}{<the closing idea>}
  \centering
  \begin{tikzpicture}[node distance=6mm and 14mm, every node/.style={font=\scriptsize}]
    \node[treenode, acc fill] (s) {<from>};
    \node[bad fill, nodebox, below left=10mm and 8mm of s] (a) {<old way>};
    \node[good fill, nodebox, below right=10mm and 8mm of s] (b) {<new way>};
    \draw[dimedge] (s) -- (a);
    \only<2->{\draw[edge, good, line width=1.2pt] (s) -- (b);}
  \end{tikzpicture}

  \begin{itemize}[<+->]
    \item <what stays the same>
    \item <what changes>
  \end{itemize}

  \only<4->{\takeaway{<the one sentence to remember>}}
\end{frame}

\end{document}
